% =====================================================================
%  Chapter 11: Hadron structure
%  Gattringer & Lang, "Quantum Chromodynamics on the Lattice" (LNP 788)
%  Converted from the raw pdftotext dump (extract/ch11.txt).
% =====================================================================
\chapter{Hadron structure}

In this section we address an important type of non-spectrum calculations.
Hadron structure is explored by matrix elements of suitable operators between hadronic states or the vacuum. The simplest such properties are the
coefficients of the hadron propagators, related to matrix elements between the
hadronic state and the vacuum, the so-called decay constants. They describe
the weak decay properties of the hadron. We discuss in some detail the lattice
calculation of decay constants as a representative of many technically similar
problems.
The matrix elements of vector or axial vector currents between single
hadron states lead to the electromagnetic and the weak form factors. Further matrix elements provide information on semileptonic decays, quark- and
gluon structure functions, and other information for effective descriptions like
the operator product expansion.
All these involve field values and as such have to be related to continuum quantities by renormalization prescriptions. Lattice perturbation theory
provides such relations. The results, however, are often not applicable in the
nonperturbative regime at the moderately small lattice spacing where one
usually works. One therefore relies also on nonperturbative determinations of
the renormalization constants that allow contact with observables given in a
continuum renormalization scheme such as $\overline{\mathrm{MS}}$.
In this chapter we will give an introductory survey of the related methods
to obtain matrix elements and renormalization factors in lattice calculations.

\section{Low-energy parameters}
In the early years of hadron theory several hypotheses were introduced to describe the observed features. The current algebra hypothesis related the (flavor
SU(3)) vector and the axial vector currents in a triplet of current commutation relations. Since the vector current is normalized due to its relation to the
electric charge, the commutation relations also lead to a normalization of the
axial vector current. Coupled to the weak interaction this current describes,
among other features, the decay of the pion. The partially conserved axial current (PCAC) hypothesis then relates the divergence of the axial vector current
to the pion field, thus defining the pion decay constant as a proportionality
factor. Another factor in the relation is the pion mass squared, indicating that
the axial charge would be conserved for exactly massless pions, i.e., Goldstone
bosons. From these assumptions many experimental results in hadron physics
can be explained. Meanwhile QCD turned out to be the underlying theory,
and what were hypotheses are now properties of that theory. In the lattice
approach we want to verify these properties and measure the corresponding
physical parameters.
In this section we consider only $u$ and $d$ quarks. These are light enough
such that the explicit breaking of chiral flavor SU(2) is small.

\subsection{Operator definitions}
In calculations of matrix elements one has to specify the relation between the
fields and the observed particle states. In quantum field theory the pion field
couples to all states with the quantum numbers $J^{PC} = 0^{-+}$ of the pion. The
pseudoscalar interpolator
\begin{equation}
P^a = \half\,\bar{\psi}\,\gamma_5\tau^a\psi
\end{equation}
has these quantum numbers. We here write the fermion fields as flavor doublets
$\psi = (u, d)$ and $\tau^a$, $a = 1, 2, 3$, are the Pauli matrices.\footnote{In our
discussion we only consider the flavor SU(2) symmetry group.} Charged quark bilinear
operators are defined through
\begin{equation}
\begin{aligned}
P^+ &\equiv P^1 - i P^2 = \half\,\bar{\psi}\gamma_5(\tau^1 - i\tau^2)\psi = \bar{d}\gamma_5 u , \\
P^- &\equiv P^1 + i P^2 = \half\,\bar{\psi}\gamma_5(\tau^1 + i\tau^2)\psi = \bar{u}\gamma_5 d .
\end{aligned}
\end{equation}

Analogously we introduce the vector and axial vector fields
\begin{equation}
\begin{aligned}
V_\mu &= \half\,\bar{\psi}\gamma_\mu\psi , &\qquad
V_\mu^a &= \half\,\bar{\psi}\gamma_\mu\tau^a\psi , &\qquad
V_\mu^\pm &= V_\mu^1 \mp i V_\mu^2 , \\
A_\mu &= \half\,\bar{\psi}\gamma_\mu\gamma_5\psi , &\qquad
A_\mu^a &= \half\,\bar{\psi}\gamma_\mu\gamma_5\tau^a\psi , &\qquad
A_\mu^\pm &= A_\mu^1 \mp i A_\mu^2 .
\end{aligned}
\end{equation}
The axial vector current is related to the vector current by commutation
relations in current algebra. For conserved vector currents the normalization
is therefore fixed for both $A_\mu$ and $V_\mu$.
One may obtain equivalent relations also directly in Euclidean space in
terms of identities of integrals over expectation values of products of field
operators [1]. In our presentation we translate the Minkowskian results of
the continuum quantum field theory to Euclidean form by replacing $it$ with
$t$. In order to facilitate switching between continuum and lattice notation we
introduce the abbreviation $\omega$ for the normalization factor of the spatial Fourier
transform, with
\begin{equation}
\omega \equiv
\begin{cases}
\dfrac{1}{(2\pi)^3} & \text{for the continuum}, \\[2mm]
|\Lambda^3| & \text{for the lattice},
\end{cases}
\end{equation}
where $|\Lambda^3| = N^3$ denotes the spatial lattice volume.
Results of a quantum field theory calculation need to be renormalized
before they can be compared to experimental numbers (see Sect.~\ref{sec:renorm}). We
denote quantities renormalized according to a continuum renormalization
scheme with a superscript $(r)$. As our reference scheme we use the
$\overline{\mathrm{MS}}$ scheme at a scale of $\mu = 2$~GeV. We define the renormalization factors $Z_P$, $Z_V$, and
$Z_A$ through
\begin{equation}
P^{(r)a} = Z_P P^a , \qquad
V^{(r)a} = Z_V V^a , \qquad
A^{(r)a} = Z_A A^a
\end{equation}
and note that they may depend on details like the lattice action and the
scale. In Sect.~\ref{sec:renorm} we discuss how one determines the renormalization factors
relating the bare lattice fields to the renormalized fields.
Also the quark mass renormalizes multiplicatively, as does the condensate
\begin{equation}
m^{(r)} = Z_m m , \qquad
\ev{\bar{\psi}\psi}^{(r)} = Z_S \ev{\bar{\psi}\psi} ,
\end{equation}
which introduces further renormalization factors $Z_m$ and $Z_S$. These are related via $Z_S = 1/Z_m$ since the combination $m\bar{\psi}\psi$ is a renormalization invariant. If chiral symmetry holds, one has the additional relations $Z_S = Z_P$ and
$Z_A = Z_V$.
If the fermionic lattice action does not obey chiral symmetry, the unrenormalized bare quark mass may have an additive contribution, called residual
mass, such that $m = m_{\mathrm{bare}} + m_{\mathrm{res}}$.
Free bosonic fields in quantum field theory obey a standard normalization.
Let us assume that the asymptotic, physical pion is described by such a renormalized triplet field with components $\phi(x)^{(r)a}$. Then in Euclidean space--time
the matrix element with the single pion state at rest $\ket{\pi^b(p = 0)}$ is
\begin{equation}
\braket{0}{\phi(x)^{(r)a}}{\pi^b(p = 0)} = \delta^{ab}\, e^{-M_\pi t} ,
\end{equation}
where $x = (\mathbf{x}, t)$. The propagator of this field reads
\begin{equation}
\braket{0}{\phi^{(r)a}(p = 0, t)\,\phi^{(r)b}(0)}{0} =
\frac{1}{2M_\pi\omega}\,\delta^{ab}\, e^{-M_\pi t} ,
\end{equation}
where $(p = 0, t)$ denotes the projection to zero spatial momentum in the
time slice $t$ (see Sect.~6.1.4). However, we still have to relate this field to the
pseudoscalar lattice field operator $P$. This will be done in the next section.
In the standard model of electroweak interaction the decay of hadrons (like
the pion or the kaon) into leptons is mediated by a -- compared to hadron
masses -- very heavy intermediate boson. The effective interaction is then
described by a Lagrangian coupling the left-handed weak to the hadronic
current, e.g., for the decay $\pi \to e\,\nu_e$ one has the coupling term
\begin{equation}
\frac{G_F}{\sqrt{2}}\,\cos\theta_c\, \bigl(\bar{u}\gamma_\mu(1-\gamma_5)d\bigr)\,
\bigl(\bar{e}\gamma_\mu(1-\gamma_5)\nu_e\bigr) .
\end{equation}
This is an example of the general operator product expansion of the weak
Hamiltonian, which will be discussed in more detail in Sect.~\ref{sec:weak}.
The pion decay according to (11.9) leads to the matrix element of the
axial vector current between the $\pi$ state and the vacuum. The relation to the
physical (renormalized) isovector pion field defines the pion decay constant
$F_\pi$, i.e.,
\begin{equation}
\partial_\mu A_\mu^{(r)a} = M_\pi^2 F_\pi\, \phi^{(r)a} .
\end{equation}
(There are also other conventions differing by, e.g., a factor of 2. Our definition corresponds to an experimental value of $F_\pi = 92.4(3)$~MeV [2]). The
divergence of the axial vector current acts as interpolating field operator of
the pion. This is an operator identity in Minkowski space; in Euclidean space
it holds in expectation values.
Various equivalent expressions for expectation values may be derived. With
(11.7) we get the so-called partially conserved axial vector current or PCAC
relation:
\begin{equation}
\partial_\mu \braket{0}{A_\mu^{(r)a}(x)}{\pi^b(p = 0)} = \delta^{ab}\, M_\pi^2 F_\pi\, e^{-M_\pi t} .
\end{equation}
These relations allow us to derive the decay constant $F_\pi$ from pair correlators between the operators $A_4$ and $\partial_t A_4$. As an example, the
$\ev{A_4 A_4}$ correlation function is dominated by the single pion at large $t$ and therefore
has the asymptotic behavior
\begin{equation}
\ev{A_4^{(r)+}(p = 0, t)\, A_4^{(r)-}(0)} \sim \frac{M_\pi F_\pi^2}{\sqrt{\omega}}\, e^{-M_\pi t} ,
\end{equation}
where the symbol $\sim$ denotes the asymptotic behavior at large $t$. The kaon
decay constant is computed analogously, with one of the quarks replaced by
the strange quark. More details will be discussed in Sect.~\ref{sec:lowenergy}.
Decay constants are but one of the quantities that may be obtained. Other
low-energy parameters include the quark masses and the condensate. For these
we introduce a new tool in the next section: Ward--Takahashi identities.

\subsection{Ward identities}
In the path integral expression for the partition function
\begin{equation}
Z = \int \mathcal{D}[\bar{\psi},\psi,U]\, e^{-S[\bar{\psi},\psi,U]}
  = \int \mathcal{D}[\bar{\psi}',\psi',U']\, e^{-S[\bar{\psi}',\psi',U']} ,
\end{equation}
the field variables $\bar{\psi},\psi, U$ are integrated over. In the second step we have used
transformed variables $\bar{\psi}',\psi', U'$ for the path integral, which in general gives the
same result. The exceptions are anomalous transformations which affect the
integration measure. This strategy of comparing the path integrations using
transformed variables was already applied successfully in Chaps.~7 and~9. Ward
identities express the invariance of the partition function (and expectation
values) under such a transformation of the field variables.
Let us study the effect of a local, unitary transformation. The expectation
value of an arbitrary interpolator $O$ is
\begin{equation}
\braket{0}{O}{0} =
\frac{1}{Z}\int \mathcal{D}[\bar{\psi},\psi,U]\, O[\bar{\psi},\psi,U]\, e^{-S[\bar{\psi},\psi,U]} .
\end{equation}
We consider an infinitesimal symmetry transformation $\psi \to \psi + \delta\psi$, $\bar{\psi} \to \bar{\psi} + \delta\bar{\psi}$ of the fermion fields in this path integral. The result has to be invariant under such a transformation of integration variables. In case of a
non-anomalous transformation the integration measure is invariant and one
gets (Ward-) identities of the form
\begin{equation}
0 = \braket{0}{\delta O}{0} - \braket{0}{O\,\delta S}{0} ,
\end{equation}
where $\delta O$ and $\delta S$ denote the linear change of the operator $O$ and the action
$S$ under the transformation. The simplest case for $O = 1$ is
\begin{equation}
\braket{0}{\delta S}{0} = 0 ,
\end{equation}
which leads to relations analogous to the classical Noether conservation laws.
For some transformations the functional interaction measure is not invariant.
This then gives rise to the so-called anomalous contributions (cf., Chaps.~7
and~9).
Let us first derive such relations in the familiar continuum form and for
$O = 1$. Later we will point out differences to the lattice notation. We study the
transformation of the action ($M = \mathrm{diag}(m_u, m_d)$ denotes the mass matrix in
flavor space)
\begin{equation}
S = \int d^4x\, \bar{\psi}(\gamma_\mu D_\mu + M)\psi , \qquad
D_\mu = \partial_\mu + iA_\mu ,
\end{equation}
with regard to infinitesimal SU(2) flavor transformations, corresponding to
the finite transformations (7.13), (7.14), (7.15), and (7.16). The linearized,
infinitesimal transformations may be written as
\begin{equation}
\begin{aligned}
\psi(x) &\to \psi'(x) = \bigl(1 + i\varepsilon(x)\lambda\bigr)\psi(x) , \\
\bar{\psi}(x) &\to \bar{\psi}'(x) = \bar{\psi}(x)\bigl(1 + i\varepsilon(x)\hat{\lambda}\bigr) ,
\end{aligned}
\end{equation}
where $\lambda$ and $\hat{\lambda}$ are products of matrices in Dirac and flavor spaces, e.g., $\lambda = 1$,
$\tau^a$, $\gamma_5$, $\gamma_5\tau^a$ and $\hat{\lambda} = -1$, $-\tau^a$, $\gamma_5$, $\gamma_5\tau^a$. We assume that $\varepsilon(x)$ is a function that
vanishes smoothly outside some bounded region. This will be of importance
since we use partial integration and the boundary terms should not contribute.
To $O(\varepsilon)$ the change of action (11.17) under such an infinitesimal transformation then is
\begin{equation}
\delta S = i \int d^4x\,
\bar{\psi}\Bigl[\varepsilon\hat{\lambda}\gamma_\mu\partial_\mu + \gamma_\mu\lambda\,\partial_\mu\varepsilon
+ i\varepsilon A_\mu(\hat{\lambda}\gamma_\mu + \gamma_\mu\lambda)
+ \varepsilon(\hat{\lambda}M + M\lambda)\Bigr]\psi .
\end{equation}
Here we study only transformations that have the property
\begin{equation}
\hat{\lambda}\gamma_\mu + \gamma_\mu\lambda = 0 ,
\end{equation}
and since the gauge field is flavor and Dirac blind, we get rid of the term
involving $A_\mu$. We have furthermore
\begin{equation}
\partial_\mu(\varepsilon\psi) = (\partial_\mu\varepsilon)\psi + \varepsilon\,\partial_\mu\psi ,
\end{equation}
leading to
\begin{equation}
\begin{aligned}
\delta S &= i \int d^4x \Bigl[\bigl(\partial_\mu\varepsilon(x)\bigr)\bar{\psi}\gamma_\mu\lambda\psi
+ \varepsilon(x)\bar{\psi}\bigl(\hat{\lambda}M + M\lambda\bigr)\psi\Bigr] \\
&= i \int d^4x\, \varepsilon(x)\Bigl[-\partial_\mu\bigl(\bar{\psi}\gamma_\mu\lambda\psi\bigr)
+ \bar{\psi}\bigl(\hat{\lambda}M + M\lambda\bigr)\psi\Bigr] .
\end{aligned}
\end{equation}
We have used integration by parts in the last step, assuming that the boundary
terms do not contribute, as guaranteed by our choice of $\varepsilon(x)$. Since this choice
is arbitrary, the condition (11.16) has to hold for any $x$. From this the operator
identity results
\begin{equation}
\partial_\mu\bigl(\bar{\psi}\gamma_\mu\lambda\psi\bigr) = \bar{\psi}\bigl(\hat{\lambda}M + M\lambda\bigr)\psi .
\end{equation}
Then, depending on the choice of $\lambda$ and $\hat{\lambda}$, but respecting (11.20), we obtain
the relations
\begin{equation}
\lambda = 1,\ \hat{\lambda} = -1 \quad\Rightarrow\quad
\partial_\mu\bigl(\bar{\psi}\gamma_\mu\psi\bigr) = 0 ,
\end{equation}
\begin{equation}
\lambda = \tau^a,\ \hat{\lambda} = -\tau^a \quad\Rightarrow\quad
\partial_\mu\bigl(\bar{\psi}\gamma_\mu\tau^a\psi\bigr) = \bar{\psi}\,[M,\tau^a]\,\psi ,
\end{equation}
\begin{equation}
\lambda = \hat{\lambda} = \gamma_5 \quad\Rightarrow\quad
\partial_\mu\bigl(\bar{\psi}\gamma_\mu\gamma_5\psi\bigr) = 2\,\bar{\psi}M\gamma_5\psi
\;(+\mathrm{anomaly}) ,
\end{equation}
\begin{equation}
\lambda = \hat{\lambda} = \gamma_5\tau^a \quad\Rightarrow\quad
\partial_\mu\bigl(\bar{\psi}\gamma_\mu\gamma_5\tau^a\psi\bigr)
= \bar{\psi}\,\{M,\tau^a\}\,\gamma_5\psi ,
\end{equation}
where we used that the mass matrix $M$ commutes with $\gamma_5$: $M\gamma_5 = \gamma_5 M$.
With the shorter notation (11.1) and (11.3) for the flavor singlet and nonsinglet vector and axial vector currents and the pseudoscalar field the relations
then assume the form
\begin{equation}
\partial_\mu V_\mu = 0 ,
\end{equation}
\begin{equation}
\partial_\mu V_\mu^a = \half\,\bar{\psi}\,[M,\tau^a]\,\psi ,
\end{equation}
\begin{equation}
\partial_\mu A_\mu = \bar{\psi}M\gamma_5\psi \;(+\mathrm{anomaly}) ,
\end{equation}
\begin{equation}
\partial_\mu A_\mu^a = \half\,\bar{\psi}\,\{M,\tau^a\}\,\gamma_5\psi .
\end{equation}

Equation (11.28) for the flavor singlet vector current is the conservation law for
the total baryon number. In the other three relations we find that the Noether
symmetry may be broken by the mass term. For mass-degenerate quarks one
has isospin symmetry and therefore $[M, \tau^a] = 0$. Consequently (11.29) defines
the conservation of the flavor vector current. In (11.30), when written for the
quantized theory, the anomaly has to be added due to a noninvariance of the
integration measure in the functional integration (cf., Sects.~7.1, 7.3, and 9.4).
Equation (11.31) is called the nonsinglet axial Ward identity (AWI): The
divergence of the flavor nonsinglet axial vector current is related to the pseudoscalar field.
We stress again that due to (11.16) these relations hold as expectation
values. Since the choice of $O$ is arbitrary we can choose the operator to be
nonvanishing and constant inside some compact, arbitrarily small space--time
region, and zero outside. Therefore (11.28), (11.29), (11.30), and (11.31) may
be considered to be valid locally. The relations are preserved under quantization and renormalization (up to the anomaly) and therefore hold for the
renormalized quantities as well.
For degenerate fermion masses $m$ the AWI (11.31) has the form
\begin{equation}
\partial_\mu A_\mu^{(r)a} = 2\,m^{(r)} P^{(r)a} .
\end{equation}
This relation is useful in many aspects. Comparing it with (11.10) and evaluating it between the vacuum and the pion we find
\begin{equation}
M_\pi^2 F_\pi \braket{0}{\phi^{(r)a}}{\pi}
= \braket{0}{\partial_\mu A_\mu^{(r)a}}{\pi}
= 2\,m^{(r)} \braket{0}{P^{(r)a}}{\pi} .
\end{equation}
This allows to determine the renormalized quark mass $m^{(r)}$ and the pion decay
constant from the asymptotic behavior of correlators of $P^{(r)a}$ and $A_\mu^{(r)a}$ (see
Sect.~\ref{sec:lowenergy}).
Further exploitation of (11.32) leads to another interesting relation. We
utilize (11.22) for transformation (11.31) by multiplying it with a local pseudoscalar operator and taking its expectation value over the Grassmann fields.
Specifically we use (11.18) with
\begin{equation}
\begin{aligned}
\lambda = \hat{\lambda} &= \gamma_5\,\half(\tau^1 + i\tau^2) , \\
O(0) &= \bar{\psi}(0)\gamma_5\,\half(\tau^1 - i\tau^2)\psi(0)
= \bar{d}(0)\gamma_5 u(0) = P^+(0) .
\end{aligned}
\end{equation}
This operator is just the pseudoscalar field. In the chiral limit $M = 0$ Eq.~(11.15) with (11.22) now leads to
\begin{equation}
-\int d^4x\, \ev{\partial_\mu A_\mu^{(r)-}(x)\, P^+(0)}
= \ev{\bar{u}(0)u(0) + \bar{d}(0)d(0)} .
\end{equation}
Inserting on the left-hand side a complete set of states, the 2-point function is saturated by the pion and the expression becomes proportional to
$\ev{P^{-}(x)P^{+}(0)}$. The left-hand side turns into
\begin{equation}
-\frac{F_\pi^2 M_\pi^4}{2\,m^{(r)}}\int d^4x\, \phi^-(x)\phi^+(0)
= -\frac{F_\pi^2 M_\pi^2}{m^{(r)}} ,
\end{equation}
where we use that the integral is twice the momentum space propagator $(p^2 +
M_\pi^2)^{-1}$ at vanishing momentum.
The right-hand side of (11.35) is the condensate $\Sigma^{(r)}$
\begin{equation}
\ev{\bar{u}(0)u(0) + \bar{d}(0)d(0)} \equiv N_f \Sigma^{(r)} ,
\end{equation}
where in our case $N_f = 2$.
We finally obtain a relation that holds for the renormalized quark mass
and the renormalized condensate $\Sigma^{(r)}$:
\begin{equation}
F_\pi^2 M_\pi^2 = -m^{(r)}\, N_f\, \Sigma^{(r)} .
\end{equation}
The pion decay constant $F_\pi$ and the pion mass $M_\pi$ do not depend on the
renormalization scheme; they are physical, measurable quantities. In our normalization the experimental value of $F_\pi$ is close to 93~MeV.
Equation (11.38) is known as Gell--Mann--Oakes--Renner (GMOR) relation
due to the seminal paper [3], where the mechanism of symmetry breaking by
the quark mass term was discussed. Note that this equation has possible
correction terms of $O(m^2)$. This is the lowest order of a systematic expansion
leading to an effective field theory known as chiral perturbation theory [4, 5].
When one attempts to derive the lattice analogies [6, 7] to the continuum
identities one runs into a problem. For actions that are not chirally symmetric,
like the Wilson action, the chiral variation of the lattice action cannot be
expressed in terms of the axial current and density alone [6] and additional
terms appear. One then argues, however, that the correlation functions of the
renormalized lattice fields should converge to the corresponding continuum
ones. This allows to derive information from expectation values like (11.58)
for nonzero distances of propagation.

\subsection{Naive currents and conserved currents on the lattice}
\label{sec:naivecurrents}
We now change to the lattice formulation and consider the construction of lattice currents corresponding to symmetry transformations of the form (11.18),
i.e.,
\begin{equation}
\delta\psi(n) = i\varepsilon(n)\lambda\psi(n) , \qquad
\delta\bar{\psi}(n) = i\bar{\psi}(n)\varepsilon(n)\hat{\lambda} .
\end{equation}
Under this transformation the action
\begin{equation}
S = \sum_{n,m} \bar{\psi}(n)\, D(U; n|m)\, \psi(m)
\end{equation}
varies, with the linear term of the variation given by
\begin{equation}
\delta S = i \sum_{n,m}
\Bigl[\bar{\psi}(n)\, D(U; n|m)\, \varepsilon(m)\lambda\, \psi(m)
+ \varepsilon(n)\,\bar{\psi}(n)\hat{\lambda}\, D(U; n|m)\, \psi(m)\Bigr] .
\end{equation}

Let us first consider Wilson fermions (5.51) and $\lambda = \tau^a$ and $\hat{\lambda} = -\tau^a$. The
term $O(\varepsilon)$ in the transformation of the Wilson action has to vanish for arbitrary $\varepsilon(n)$. After shifting the index for the summation over the lattice sites
this can be cast into a local form. We end up with an equation equivalent
to (11.29), with the replacement of the continuum derivatives by a nearest
neighbor difference
\begin{equation}
\sum_{\mu=1}^{4} \frac{1}{a}\Bigl[V_\mu^a(n + \hat{\mu}) - V_\mu^a(n)\Bigr]
= \half\,\bar{\psi}(n)\,[M, \tau^a]\,\psi(n) ,
\end{equation}
with the lattice vector current
\begin{equation}
V_\mu^a(n) = \half \Bigl[
\bar{\psi}(n + \hat{\mu})(1 + \gamma_\mu)U_\mu(n)^\dagger\tau^a\psi(n)
- \bar{\psi}(n)(1 - \gamma_\mu)U_\mu(n)\tau^a\psi(n + \hat{\mu})
\Bigr] .
\end{equation}

We can write (11.42) as
\begin{equation}
\Delta_\mu V_\mu^a(n) = \half\,\bar{\psi}(n)\,[M, \tau^a]\,\psi(n) ,
\end{equation}
where $\Delta_\mu$ is the forward lattice derivative $\Delta_\mu f(n) \equiv \bigl(f(n + \hat{\mu}) - f(n)\bigr)/a$.
This agrees (up to higher orders in the lattice spacing) with the corresponding
continuum expression (11.29). For $u$ and $d$ quarks with identical masses the
current is conserved like in the continuum.
For the axial transformations, however, this simple equivalence does not
work any more for fermion actions which violate chiral symmetry (like the
Wilson action). There are symmetry-breaking terms divergent with $1/a$ and it
is not possible to define an exactly conserved axial vector current for vanishing
quark masses.
In Chap.~7 we discussed Dirac operators which obey the Ginsparg--Wilson
relation and thus chiral symmetry on the lattice. The corresponding operators, e.g., the overlap operator, are more complicated. It therefore pays off
to consider the problem of constructing Noether currents for a general action
without specifying a particular form explicitly. For this we follow [8] and discuss as an example a simple $n$-dependent scalar transformation like in (11.18)
and (11.24) with $\lambda = 1$, $\hat{\lambda} = -1$.
The variation (11.41) assumes the form ($k$ is a position index here)
\begin{equation}
\begin{aligned}
\delta S &= i a \sum_{k,\mu} \bigl(\Delta_\mu\varepsilon(k)\bigr) J_\mu(k)
= i \sum_{k,\mu} \bigl(\varepsilon(k + \hat{\mu}) - \varepsilon(k)\bigr) J_\mu(k) \\
&= i \sum_{k,\mu} \varepsilon(k) \bigl(J_\mu(k - \hat{\mu}) - J_\mu(k)\bigr)
= -i a \sum_{k,\mu} \varepsilon(k)\, \Delta_\mu^* J_\mu(k) ,
\end{aligned}
\end{equation}
where $\Delta_\mu$ and $\Delta_\mu^*$ denote forward and backward lattice derivatives. In case
the action is invariant under such a change of variables, i.e., $\delta S = 0$, we could
identify $J_\mu(k)$ with the conserved Noether current on the lattice.

For this end let us consider a transformation of each individual link
variable:
\begin{equation}
U_\mu'(n) = U_\mu^{(\alpha)}(n) \equiv e^{i\alpha_\mu(n)} U_\mu(n)
\equiv e^{i\varepsilon(n)} U_\mu(n)\, e^{-i\varepsilon(n + \hat{\mu})} ,
\end{equation}
where $\alpha_\mu(n) = -\Delta_\mu\varepsilon(n)$. We have rewritten this transformation in order
to utilize the transformation properties of the Dirac operator under gauge
transformations, hence
\begin{equation}
\begin{aligned}
&U_\mu(n) \to U_\mu'(n) \qquad\Rightarrow\qquad\\
&D(U; n|m) \to D(U'; n|m) = e^{i\varepsilon(n)}\, D(U; n|m)\, e^{-i\varepsilon(m)} .
\end{aligned}
\end{equation}

Taking the difference between transformed and untransformed Dirac operators
we find at $O(\varepsilon)$:
\begin{equation}
D(U'; n|m) - D(U; n|m)
= i\bigl(D(U; n|m)\,\varepsilon(m) - \varepsilon(n)\, D(U; n|m)\bigr) ,
\end{equation}
where we have taken into account just the leading order in $\varepsilon$. This, however,
is the kernel of (11.41). Alternatively, we can determine this leading order
difference also by a derivative with regard to $\alpha_\mu$:
\begin{equation}
\begin{aligned}
D(U'; n|m) - D(U; n|m)
&= -\sum_{k,\mu} \alpha_\mu(k)
\left.\left(\frac{\partial D\bigl(U_\mu^{(\alpha)}(k); n|m\bigr)}
{\partial\alpha_\mu(k)}\right)\right|_{\alpha=0} \\
&= \sum_{k,\mu} \bigl(\varepsilon(k + \hat{\mu}) - \varepsilon(k)\bigr)
\left.\left(\frac{\partial D\bigl(U_\mu^{(\alpha)}(k); n|m\bigr)}
{\partial\alpha_\mu(k)}\right)\right|_{\alpha=0} .
\end{aligned}
\end{equation}

Comparison with (11.45) leads to the identification of the Noether current
\begin{equation}
\begin{aligned}
J_\mu(k) &= \sum_{n,m} \bar{\psi}(n)\, K_\mu(k; n|m)\, \psi(m) , \\
K_\mu(k; n|m) &\equiv -i
\left.\left(\frac{\partial D\bigl(U_\mu^{(\alpha)}(k); n|m\bigr)}
{\partial\alpha_\mu(k)}\right)\right|_{\alpha=0} .
\end{aligned}
\end{equation}
We find that if we imagine the Dirac operator as a sum over paths connecting
$n$ with $m$, then all paths which run through the link $U_\mu(k)$ contribute to
$J_\mu(k)$. From the construction the gauge invariance of $J_\mu(k)$ is obvious. In
practical applications the derivatives have to be performed numerically by
taking differences. Applying this construction to the Wilson Dirac operator
we recover the simple local currents as, e.g., in (11.43).
For chirally symmetric actions it is possible to construct conserved and
covariantly transforming vector currents as well as covariant scalar densities.
These then obey Ward identities like in the continuum.
Following the strategy outlined in our example, in [8, 9] chiral covariant
and conserved currents have been constructed for general, exactly
chiral-invariant actions like the overlap action. We refer to the original papers
for the explicit derivation and only quote the results. With the abbreviation
\begin{equation}
\hat{\gamma}_5 = \gamma_5 (1 - D) ,
\end{equation}
we have
\begin{equation}
\begin{aligned}
V_\mu^a(k) &=
\frac{1}{4}\sum_{n,m} \bar{\psi}(n)\Bigl[K_\mu(k; n|m)
- \gamma_5\sum_j K_\mu(k; n|j)\,(\hat{\gamma}_5)_{jm}\Bigr]\tau^a\psi(m) , \\[2mm]
A_\mu^a(k) &=
\frac{1}{4}\sum_{n,m} \bar{\psi}(n)\Bigl[-\gamma_5 K_\mu(k; n|m)
+ \sum_j K_\mu(k; n|j)\,(\hat{\gamma}_5)_{jm}\Bigr]\tau^a\psi(m) , \\[2mm]
S^a(k) &=
\frac{1}{4}\sum_{n,m} \bar{\psi}(n)\Bigl[\delta_{nk}\bigl(1 - \half D\bigr)_{km}\\
&\qquad + \half\sum_j \bigl(1 - \half D\bigr)_{nj}\delta_{jk}\delta_{mk}
+ (\hat{\gamma}_5)_{jk}(\hat{\gamma}_5)_{km}\Bigr]\tau^a\psi(m) , \\[2mm]
P^a(k) &=
\frac{1}{4}\sum_{n,m} \bar{\psi}(n)\gamma_5\Bigl[\delta_{nk}\bigl(1 - \half D\bigr)_{km}\\
&\qquad + \half\sum_j \bigl(1 - \half D\bigr)_{nj}(\hat{\gamma}_5)_{jk}\delta_{mk}
+ \delta_{jk}(\hat{\gamma}_5)_{km}\Bigr]\tau^a\psi(m) .
\end{aligned}
\end{equation}
These conserved currents and equivalently constructed scalar and pseudoscalar
densities transform covariantly under the chiral transformations, like in the
continuum. For conserved vector currents the renormalization is trivial (the
charge is conserved and the renormalization constants obey $Z_A = Z_V = 1$).
An important advantage is that the currents and densities (dimension-3 operators) are automatically $O(a)$ improved and do not mix with dimension-4
or dimension-3 operators, respectively.
Since the implementation of prescription (11.50) is rather involved, the
simpler expressions
\begin{equation}
\begin{aligned}
\bar{V}_\mu^a(k) &=
\half \sum_m \bar{\psi}(k)\gamma_\mu\bigl(1 - \half D\bigr)_{km}\tau^a\psi(m) , \\[1mm]
\bar{A}_\mu^a(k) &=
\half \sum_m \bar{\psi}(k)\gamma_\mu\gamma_5\bigl(1 - \half D\bigr)_{km}\tau^a\psi(m) , \\[1mm]
\bar{S}^a(k) &=
\half \sum_m \bar{\psi}(k)\bigl(1 - \half D\bigr)_{km}\tau^a\psi(m) , \\[1mm]
\bar{P}^a(k) &=
\half \sum_m \bar{\psi}(k)\gamma_5\bigl(1 - \half D\bigr)_{km}\tau^a\psi(m)
\end{aligned}
\end{equation}
have been used. They are also covariant but not conserved and thus not related
by Ward identities [8]. For direct constructions based on Ginsparg--Wilson
fermions see also [10, 11].
In correlation functions over distances larger than the range of the Dirac
operator the covariant densities in (11.52) can be replaced by the even simpler,
point-like expressions
\begin{equation}
\begin{aligned}
S^a(k) &\approx \half\, \frac{1}{1 - m/2}\,\bar{\psi}(k)\tau^a\psi(k) , \\
P^a(k) &\approx \half\, \frac{1}{(1 - m/2)(1 - m^2/4)}\,\bar{\psi}(k)\gamma_5\tau^a\psi(k) .
\end{aligned}
\end{equation}
This comes about since the propagator involves the inverse massive Dirac
operator $D_m^{-1}$ and from (7.41) and (7.42) we know that
\begin{equation}
1 - \half D = \frac{1}{1 - m/2}\Bigl(1 - \half D_m\Bigr) .
\end{equation}
The insertions $1 - \half D$ in (11.53) thus can be expressed in terms of the massive
Dirac operator $D_m$. When considering correlation functions of the interpolators (11.53) the insertions give rise to products $D_m D_m^{-1}$ which just contribute
locally (``contact term'') and may be omitted [8].
In the continuum limit the currents renormalize multiplicatively with factors like $Z_V$ and $Z_A$ and we will discuss the determination of these factors in
Sect.~\ref{sec:renorm}.

\subsection{Low-energy parameters from correlation functions}
\label{sec:lowenergy}
In the last section we have shown that the continuum identities for correlation
functions may be translated to the lattice, although the identification of the
currents and densities on the lattice is not straightforward. We can expect
that the asymptotic behavior of correlation functions may be compared. Thus
we may utilize equations like (11.12) to determine low-energy constants from
the asymptotic (long time distance) behavior of certain correlation functions
and ratios thereof.
The renormalized quark mass may be determined utilizing the axial Ward
identities (11.31), (11.33) and considering the (large $t$) asymptotic behavior of
\begin{equation}
\frac{1}{2}\,
\frac{\ev{0}{\partial_t A_4^{-(r)}(p = 0, t)\,X(0)}{0}}
{\ev{0}{P^{-(r)}(p = 0, t)\,X(0)}{0}}
= \frac{1}{2}\,
\frac{Z_A \ev{0}{\partial_t A_4^-(p = 0, t)\,X(0)}{0}}
{Z_P \ev{0}{P^-(p = 0, t)\,X(0)}{0}}
\sim m^{(r)} .
\end{equation}
For $X$ one chooses an interpolator coupling to the pion and often one takes
$X = P$. Numerator and denominator are both asymptotically dominated by
the pion intermediate state and the asymptotic $t$-dependence cancels. In the
calculations one can identify a plateau as soon as the asymptotic behavior has
set in. In this relation one needs to know the renormalization constants $Z_A$
and $Z_P$ relating the lattice interpolators with the continuum normalization.
These can be determined separately, as will be discussed later.
Neglecting the renormalization constants one often just determines the
asymptotic ratio of the unrenormalized correlators:
\begin{equation}
\frac{1}{2}\,
\frac{\ev{0}{\partial_t A_4^-(p = 0, t)\,X(0)}{0}}
{\ev{0}{P^-(p = 0, t)\,X(0)}{0}}
\sim m_{\mathrm{AWI}} .
\end{equation}
Due to its origin this ratio is called the AWI mass $m_{\mathrm{AWI}}$ or PCAC mass. We
have already used (11.57) as a tool for determining the coefficients $c_{sw}$ and $c_A$
in the Symanzik improvement program in Sect.~9.1.4.
The quark mass of the lattice action may undergo two renormalizations.
In case the action is not chirally symmetric there may be an additive renormalization called residual mass $m_{\mathrm{res}}$ and $m_{\mathrm{AWI}} = m_{\mathrm{bare}} + m_{\mathrm{res}}$. We then have
$m^{(r)} = (Z_A / Z_P)\, m_{\mathrm{AWI}}$. In relation to, e.g., the continuum $\overline{\mathrm{MS}}$ scheme one has
a multiplicative renormalization factor $m^{(r)} = Z_m m$. For chirally symmetric
actions $Z_m Z_P = 1$ and thus the bare mass in lattice units is $m = Z_A\, m_{\mathrm{AWI}}$.
The AWI mass vanishes at the points in the space of couplings, where
chiral symmetry holds. This is particularly useful in cases with additive mass
renormalization, since it allows to identify the residual mass (see below). For
the Wilson action one identifies the critical value of the hopping parameter
$\kappa_c(\beta)$ with its help: It is defined as the value where $m_{\mathrm{AWI}}$ vanishes and the
pions are massless (see also Sect.~9.1.4).
The pion decay constant may be obtained, e.g., from the axial current
correlation function (11.12):
\begin{equation}
Z_A^2 \ev{A_4^+(p = 0, t)\, A_4^-(0)} \sim \frac{M_\pi F_\pi^2}{\sqrt{\omega}}\, e^{-M_\pi t} .
\end{equation}
Due to (11.33) we also find
\begin{equation}
\ev{0}{P^+(p = 0, t)\,P^-(0)}{0} \sim \frac{A}{\sqrt{\omega}}\, e^{-M_\pi t}
\qquad\text{with}\qquad
A = \frac{M_\pi^3 F_\pi^2}{4\,m^{(r)2} Z_P^2} .
\end{equation}
From the coefficient to the asymptotic behavior we get
\begin{equation}
F_\pi = 2\,Z_P\, m^{(r)}\, \sqrt{\frac{A}{M_\pi^3}}
= 2\,Z_A\, m_{\mathrm{AWI}}\, \sqrt{\frac{A}{M_\pi^3}} .
\end{equation}
The GMOR equation (11.38) may be used to compute the quark condensate, once $F_\pi$, $M_\pi$, and $m^{(r)}$ are known. One may also use correlation functions
like
\begin{equation}
Z_A Z_P\, \bigl|\ev{0}{A_4^+(p = 0, t)\,P^-(0)}{0}\bigr| \sim \frac{|\Sigma^{(r)}|}{\sqrt{\omega}}\, e^{-M_\pi t}
\end{equation}
for its determination (cf. [12] for more such relations).

\section{Renormalization}
\label{sec:renorm}
\subsection{Why do we need renormalization?}
When quantizing a field theory one has to regularize it, and the lattice formulation is one possible ultraviolet regulator. Identification of computed quantities with physical observables then completes the renormalization process.
Different regulators and actions require different renormalization parameters.
Some observables like hadron masses can be measured directly as dimensionless numbers $a\,m_{\mathrm{phys}}$ from the exponential decay of the correlation functions. By comparing one such number with the physical mass $m_{\mathrm{phys}}$ of that
particle one may determine the lattice spacing $a$ and establish the scale (see
Sect.~6.3). Other quantities like decay constants or form factors involve fields
and their relationship to experiments and continuum normalization is more
involved.
Matrix elements and low-energy parameters like quark masses or the condensate play an important role in QCD phenomenology. However, their values
depend on the definitions in some renormalization scheme. Since one wants
to compare with physical, experimentally measured quantities one has to relate these to the parameters of the underlying formulation of the quantized
theory. For that purpose one has to determine the scaling properties, i.e., the
dependence on the scale $a$, of the observables and their behavior under renormalization. Many quantities are intrinsically scale dependent and are even
divergent when the cutoff parameter $a$ is removed. We can distinguish the
following scenarios:
\begin{itemize}
\item Finite operators, e.g., vector currents and axial vector currents or the ratios
of scalar and pseudoscalar densities: The multiplicative renormalization
factors $Z_V$, $Z_A$ and the ratio $Z_S/Z_P$ assume finite values for vanishing
regularization parameter $a$, thus they should be scale independent. For
conserved currents and chirally symmetric actions $Z_V = Z_A = 1$. Depending on the Dirac operator, it may be complicated to use the conserved
currents (see Sect.~\ref{sec:naivecurrents}) and one therefore relies on simpler, e.g., pointlike currents where the $Z$-factors are not known a priori.
\item Logarithmically divergent operators (which are, e.g., scalar and pseudoscalar densities or flavor $\Delta F = 2$ changing 4-fermion operators) are
scale dependent.
\item Power divergent operators: These occur in regularizations with intrinsic
mass scale like in the lattice regularizations. Then mixing with lower
dimensional operators becomes possible, leading to power divergences
$\order{1/a^n}$ [13]. Examples are $\order{1/m}$ terms in heavy quark effective theory (HQET) discussed in Sect.~10.4.
\end{itemize}

In the regularization some symmetries of the original theory are lost. In
the lattice approach these are, e.g., the continuous space--time transformations
and (for some actions) the chiral symmetry. Exact GW fermions are protected
by their chiral symmetry and renormalization of operators constructed from
them is simpler than for, e.g., Wilson fermions. Chiral symmetry implies several relations between renormalization constants, e.g., $Z_A = Z_V$ and $Z_S = Z_P$.
In order to compute renormalization constants nonperturbatively on the
lattice one needs a renormalization scheme which can be implemented in lattice Monte Carlo simulations and in continuum perturbation theory. The latter
property is necessary to enable the conversion of the lattice results to the continuum scheme. The most favored continuum scheme is the modified minimal
subtraction $\overline{\mathrm{MS}}$ scheme. In such a mass-independent scheme the renormalization factors depend only on a normalization mass $\mu$ and the coupling constant,
but not on quark masses.

\subsection{Renormalization with the Rome--Southampton method}
Since one compares the lattice results with results of the continuum one could
utilize lattice perturbation theory (see [13] and references therein). This has
been done but in many cases the convergence is too bad to lead to reliable
numbers. One therefore often relies on nonperturbative methods [14].
The idea is to compare bare lattice correlation functions, couplings, and
masses determined in nonperturbative lattice calculations with quantities
in the so-called Regularization-independent (RI) scheme (often also called
RI/MOM scheme). The connection between the RI quantities and those defined in the $\overline{\mathrm{MS}}$ scheme can then be done in continuum perturbation theory
[15].
As an example we discuss the renormalization factors of quark bilinear
operators. To be more explicit, let us consider the local, flavor nonsinglet
quark field bilinear operators
\begin{equation}
O_\Gamma \equiv \bar{u}\,\Gamma\, d .
\end{equation}
Here $\Gamma$ denotes a Clifford algebra matrix. According to their Lorentz symmetry we denote the five types of $\Gamma$ by $S$, $V$, $A$, $T$, and $P$ corresponding to
scalar, vector, axial vector, tensor, and pseudoscalar (corresponding to $1$, $\gamma_\mu$,
$\gamma_\mu\gamma_5$, $\frac{i}{2}[\gamma_\mu, \gamma_\nu]$, and $\gamma_5$).
One studies expectation values $\braket{p}{O_\Gamma}{p}$ (i.e., at zero momentum transfer) of the bilinear quark operators between quark fields $\ket{p}$ at a specific momentum value $p^2 = \mu^2$ and matches them to the corresponding tree-level
matrix element:
\begin{equation}
Z_\Gamma\, \braket{p}{O_\Gamma}{p}\big|_{p^2 = \mu^2} = \braket{p}{O_\Gamma}{p}_0\big|_{p^2 = \mu^2} .
\end{equation}
The renormalization constant $Z_\Gamma$ is the proportionality factor between the
interacting and the free case.
This procedure is expected to work in a window
\begin{equation}
\Lambda_{\mathrm{QCD}}^2 \ll \mu^2 \ll 1/a^2 ,
\end{equation}
where discretization effects can be neglected, because the renormalization
scale $\mu$ is small compared with the lattice cutoff $1/a$. Then (few-loop) continuum perturbation theory can be used to connect different schemes, because $\mu$
is much larger than the QCD scale parameter $\Lambda_{\mathrm{QCD}}$. For comparing with the
$\overline{\mathrm{MS}}$ scheme a typical value is $\mu = 2$~GeV. In most calculations one has $(a\mu) \approx 1$
(or even somewhat larger) and the upper limit is not strictly obeyed. On the
other hand, the limit also depends on the scaling properties of the actions
involved.
Since (11.63) is gauge variant, one has to work in a fixed gauge and must
compare the gauge dependent lattice matrix elements with the continuum
results in the same gauge. Landau gauge fixing is a suitable choice, but one has
to keep in mind that the Gribov copies uncertainty could spoil the comparison.
In the lattice calculations one finds little, if any, signal of such an effect [16--18].
We summarize the method following [14] in the modification of [19]. Note
that in (11.63) one compares matrices in color and Dirac spaces. Taking the
trace one obtains for the renormalization condition
\begin{equation}
\frac{Z_\Gamma}{12}\,
\tr\Bigl[\braket{p}{O_\Gamma}{p}\,\bigl(\braket{p}{O_\Gamma}{p}_0\bigr)^{-1}\Bigr]
\Big|_{p^2 = \mu^2} = 1 .
\end{equation}
The matrix element
\begin{equation}
\braket{p}{O_\Gamma}{p} = \frac{1}{Z_q}\,\Lambda_\Gamma(p)
\end{equation}
is proportional to the amputated Green function
\begin{equation}
\Lambda_\Gamma(p) = S^{-1}(p)\, G_\Gamma(p)\, S^{-1}(p) ,
\end{equation}
and $Z_q$ is the quark field renormalization constant to be discussed below. The
Green function $G_\Gamma(p)$ is determined as the expectation value:
\begin{equation}
G_\Gamma(p)_{\alpha\beta}^{ab} =
\frac{1}{V}\sum_{x,y,z} e^{-ip(x-y)}
\ev{u_\alpha^a(x)\, O_\Gamma(z)\, \bar{d}_\beta^b(y)} .
\end{equation}
The indices $\alpha, \beta$ and $a, b$ run over Dirac and color indices, respectively, and
$V$ denotes the lattice volume. The quark propagator is
\begin{equation}
S_{\alpha\beta}^{ab}(x, y) = \ev{u_\alpha^a(x)\, \bar{u}_\beta^b(y)}
= \ev{d_\alpha^a(x)\, \bar{d}_\beta^b(y)}
\end{equation}
(assuming that $u$ and $d$ have equal masses and using the Landau gauge for
the expectation value).
So we have to compute $G_\Gamma(p)$ and $S(p)$. This is done in the following way.
For the quark propagator $S^{(n)}$ evaluated on a single gauge configuration $n$ we
define
\begin{equation}
S^{(n)}(x|p) = \sum_y e^{ipy}\, S^{(n)}(x, y) .
\end{equation}
Taking into account $\gamma_5$-hermiticity of the propagator we may, for quark bilinear operators $O_\Gamma$ as defined in (11.62), rewrite $G_\Gamma(p)$ in terms of the quantities
(11.70):
\begin{equation}
G_\Gamma(p) =
\frac{1}{V}\sum_{x,y,z} e^{-ip(x-y)}
\ev{u(x)\,\bar{u}(z)\,\Gamma\, d(z)\,\bar{d}(y)} ,
\end{equation}
\begin{equation}
G_\Gamma(p) \approx \frac{1}{VN}\sum_{n=1}^{N}\sum_z
\gamma_5\, S^{(n)}(z|p)^\dagger\, \gamma_5\, \Gamma\, S^{(n)}(z|p) ,
\end{equation}
where we approximate the expectation value by averaging over $N$ gauge configurations. Similarly we find the quark propagator in momentum space
\begin{equation}
S(p) \approx \frac{1}{VN}\sum_{n=1}^{N}\sum_x e^{-ipx}\, S^{(n)}(x|p) .
\end{equation}
$S^{(n)}(y|p)$ is computed by solving the lattice Dirac equation ($D$ denotes the
Dirac operator)
\begin{equation}
\sum_y D(z, y)\, S^{(n)}(y|p) = e^{ipz}
\end{equation}
with a momentum source (cf. [19]). This has the disadvantage that one has to
determine the quark propagators for several momentum sources, whereas in
the original method [14] one uses point sources (i.e., taking into account just
$z = 0$ instead of summing over all $z$) and projects the quark sink to the desired
momentum values. However, using momentum sources has the big advantage
of a significantly better signal.
The quark field renormalization constant is obtained by comparing the
quark propagator to the free (lattice) propagator. Using the so-called RI$'$
scheme we take
\begin{equation}
Z_q' =
\frac{1}{12}\,
\tr\left[S^{-1}(p)\,
\frac{R(p) - i\gamma_\nu v_\nu(p)}
{R(p)^2 + \sum_\nu v_\nu(p)^2}\right]
\Bigg|_{p^2 = \mu^2} .
\end{equation}
The RI$'$ scheme differs from the RI scheme only by the definition of the quark
field renormalization constant. In (11.75), $R$ and $v_\nu$ are the scalar and vector
terms appearing in the free Dirac operator, which in momentum space reads
$D(p) = i\gamma_\nu v_\nu(p) + R(p)$. Using Fourier transformation one can compute $R$
and $v_\nu$ from the definition of $D$. They are normalized such that one finds
\begin{equation}
v_\nu(p) = i p_\nu + \order{ap^2} , \qquad
R(p) = \order{ap^2} .
\end{equation}
Landau gauge fixing $\partial_\nu A_\nu = 0$ is implemented as discussed in [17, 20] by iteratively minimizing a functional of the link variables with stochastic overrelaxation [21] (cf., Sect.~3.2.2). As is well known this type of gauge fixing still
allows for Gribov copies corresponding to further local minima. This gives
rise to an uncertainty which one has to check, e.g., by studying the effect of
random gauge transformations.
It is easy to check that at tree level one finds $\braket{p}{O_\Gamma}{p}_0 = \Gamma$. Putting
things together we obtain the final formula for $Z_\Gamma$ in the RI$'$ scheme:
\begin{equation}
Z_\Gamma^{\mathrm{RI}'} = \frac{12\,Z_q'}{\tr\bigl[\Lambda_\Gamma(p)\,\Gamma^{-1}\bigr]}
\Bigg|_{p^2 = \mu^2} .
\end{equation}
For $\Gamma = \gamma_\nu, \gamma_\nu\gamma_5$ averaging over the index $\nu$ under the trace is implied.
The connection between different renormalization schemes is established
using continuum perturbation theory. In the two schemes $\overline{\mathrm{MS}}$ and RI$'$ the
renormalization factors are related, and we obtain the final result for the
lattice numbers $Z_\Gamma^{\mathrm{RI}'}$ as
\begin{equation}
Z_\Gamma^{\overline{\mathrm{MS}}}(\mu^2) = R_\Gamma(\mu^2)\, Z_\Gamma^{\mathrm{RI}'}(\mu^2) ,
\end{equation}
where the ratio $R_\Gamma(\mu^2)$ may be computed in perturbation theory. Both,
$Z_\Gamma^{\mathrm{RI}'}(\mu^2)$ and $Z_\Gamma^{\overline{\mathrm{MS}}}(\mu^2)$ may be divergent in the continuum limit, the ratio
stays finite, though.
Conversion factors between RI, RI$'$, and $\overline{\mathrm{MS}}$ schemes have been determined in Landau gauge and 3-loop order in [15, 22, 23]. Integrating the renormalization group differential equations one can also define scale-independent
quantities $Z^{\mathrm{RGI}}$ (RGI = renormalization group invariant) [19, 24].
Here we have only discussed the simple case of renormalizing quark bilinears of the form (11.62). Operators with derivatives, as they are used for
matrix elements discussed in Sect.~\ref{sec:matrixelements}, may be treated with similar methods.
In the RI scheme one has to extrapolate to the chiral limit for comparison with the mass-independent $\overline{\mathrm{MS}}$ scheme. An alternative approach is the
Schr\"odinger functional method (see Sect.~9.1), where one may work directly
at the chiral point [13, 25].

\section{Hadronic decays and scattering}
\subsection{Threshold region}
In finite volumes the observables show finite size effects. Let us first discuss the
propagation of a single hadron. We assume that the time extent is sufficiently
large and thus not important for finite size effects. We also assume periodic
boundary conditions for the spatial directions of the volume. Then there are
two ways the single particle propagator is affected by the finite volume: On
the one hand the propagator between two points separated by some distance
in Euclidean time is actually a sum over propagators between the source and
all spatial mirror images. On the other hand mass renormalizing interactions
(like loops) will also have paths running around a periodic spatial direction
before coming back to the vertex, as in Fig.~11.1. One could interpret this as
a squeezing effect of the finite box on the virtual polarization cloud of the
propagating particle. The leading effect comes from the lightest interacting
particles in the system, i.e., the pions in QCD. The control parameter is the
ratio of the box size $L$ to the light particle correlation length, or, equivalently,
the product $M_\pi L$.

\begin{figure}[htbp]
  \centering
  \includegraphics[width=0.8\textwidth]{fig111.png}
  \caption*{Fig. 11.1. \emph{Propagation of a pion. Left: free (above) and with
  interaction loop (below). Right: additional terms in a finite, periodic spatial
  box, where the pion also can run around the torus due to periodic boundary
  conditions}.}
  \label{fig:111}
\end{figure}

This has been discussed in [26] for the continuum quantum field theory of
hadrons. In QCD we have quarks and gluons as the basic objects. However,
due to confinement the quark loops are suppressed and the leading effect is
expected to be due to pions. Thus the derivation should be applicable to QCD.
A single hadron does not change its momentum when emitting and absorbing a pion, as in Fig.~11.1. Thus the interaction is related to the forward
scattering amplitude. If, in an approximation, this is just a constant $\alpha$, then
the mass change is proportional to the (divergent) loop contribution
\begin{equation}
m^2 = m_0^2 + \alpha \int \frac{d^4p}{(2\pi)^4}\, G(p) = m_0^2 + \alpha\, G(0) ,
\end{equation}
where $G(p)$ is the propagator in momentum space. The integral gives the real
space propagator at the origin $G(0)$, as made explicit in the second step of
(11.79). The divergence is absorbed in the definition of the bare mass (renormalization).
We now consider a finite, periodic, spatial box of size $L^3$ and temporal
extent large enough to neglect modifications due to its finiteness. In such a
finite volume we have to replace the real space propagator by a sum
\begin{equation}
G(x) \Rightarrow \sum_{s \in \mathbb{Z}^3} G(x + L s) ,
\end{equation}
where $s \in \mathbb{Z}^3$ gives the shift between the original lattice and its mirror images.
The mass difference between finite and infinite box due to the interaction
term, i.e., due to the contributions from propagation from the mirror images,
is given by
\begin{equation}
m^2(L) - m^2 = \alpha \sum_{s \neq 0} G(L s) .
\end{equation}
In order to estimate these contributions we now need the large distance behavior of the free boson propagator. There are various ways to determine the
Fourier transform
\begin{equation}
G(x) = \int \frac{d^4p}{(2\pi)^4}\, \frac{e^{ipx}}{p^2 + m^2}
\sim \frac{m^2 \sqrt{8\pi}\, e^{-m|x|}}{(4\pi)^2 (m|x|)^{3/2}} .
\end{equation}
We therefore can expect a leading size dependence $\propto e^{-mL}/L^{3/2}$. The
coefficient can be determined, e.g., from ChPT [27--29].
L\"uscher [26] sums all orders of perturbation theory of the loop contributions on a periodic, finite lattice and obtains the expression
\begin{equation}
m(L) - m = -\frac{1}{2}\,\frac{3}{(4\pi)^2 mL}
\int_{-\infty}^{\infty} dy\, e^{-\sqrt{y^2 + m^2}\, L}\, F(iy) + O(e^{-mL}) ,
\end{equation}
where the second term is suppressed since $m > \tfrac{3}{2}m$ can be shown. $F(iy)$
denotes the $X\pi \to X\pi$ forward-scattering amplitude $F(\nu = s - u) = A(s, t =
0, u)$, analytically continued to imaginary argument, where $s$, $t$, and $u$ are the
relativistically invariant Mandelstam variables.
We will not reproduce the derivation here but just make the result plausible. For this we integrate (11.82) by choosing $x$ along the time axis and write
\begin{equation}
\begin{aligned}
G(x) &= \int_{-\infty}^{\infty} \frac{dE}{2\pi}\, e^{iEx}
\int_{-\infty}^{\infty} \frac{d^3p}{(2\pi)^3}\, \frac{1}{E^2 + \mathbf{p}^2 + m^2} \\
&= \int_{-\infty}^{\infty} \frac{dE}{2\pi}\, e^{iEx}
\int_0^{\infty} \frac{dy}{2\pi}\, \frac{y^2}{E^2 + y^2 + m^2}
\end{aligned}
\end{equation}
(here $y = |\mathbf{p}|$). The integral over $E$ may be solved using Cauchy's theorem,
i.e., completing the contour by a half-circle at infinity and picking up the pole
at $E = i\sqrt{y^2 + m^2}$ in the upper half-plane,
\begin{equation}
\begin{aligned}
G(x) &= \frac{1}{4\pi^2} \int_0^{\infty} dy\,
\frac{y^2}{\sqrt{y^2 + m^2}}\, e^{-\sqrt{y^2 + m^2}\, x} \\
&= -\frac{1}{4\pi^2}\,\frac{1}{x} \int_0^{\infty} dp\, p\,
\frac{\partial}{\partial p}\, e^{-\sqrt{p^2 + m^2}\, x}
= \frac{1}{8\pi^2 x} \int_{-\infty}^{\infty} dp\, e^{-\sqrt{p^2 + m^2}\, x} .
\end{aligned}
\end{equation}
In the second step we did a partial integration and then doubled the integration interval to get a form resembling (11.83). We now add up just the leading
contributions of the sum (11.81) due to the six neighbor images and get
\begin{equation}
m^2(L) - m^2 = \frac{6\alpha}{8\pi^2 L}
\int_{-\infty}^{\infty} dp\, e^{-\sqrt{p^2 + m^2}\, L} ,
\end{equation}
and from this
\begin{equation}
m(L) - m \approx \frac{m^2(L) - m^2}{2m}
= \frac{3\,(-2\alpha)}{2m}\,
\frac{1}{(4\pi)^2 L}
\int_{-\infty}^{\infty} dp\, e^{-\sqrt{p^2 + m^2}\, L} .
\end{equation}
This is just (11.82) with the scattering amplitude replaced by the constant
coupling $F \to -2\alpha$. Thus we reproduce the simple leading behavior discussed
earlier.
In the second paper of the series [30] L\"uscher then relates the energy
spectrum of two stable particles (at rest) in a finite, periodic box to their
elastic scattering amplitude. Whereas the finite volume effects for the single
particle are exponentially small, now the interaction between the two particles
leads to only polynomial suppression $O(1/L^3)$. The physical picture is that the
localized particles need to be close to each other to interact and the probability
to find one particle close to another is inversely proportional to the spatial
volume.
Near threshold the scattering amplitude is dominated by the $s$-wave scattering length $a_0 = \lim_{k\to 0} \delta_0(k)/k$ (where $k$ is the 2-particle momentum and
$\delta_0$ the $s$-wave scattering phase shift). For a 2-particle system the expansion of
the ground state energy in $1/L$ is
\begin{equation}
W = 2M_\pi - \frac{4\pi a_0}{M_\pi L^3}
\left[1 + c_1\frac{a_0}{L} + c_2\frac{a_0^2}{L^2} + O\left(\frac{1}{L^3}\right)\right] .
\end{equation}
The coefficients $c_1 = -2.837297$ and $c_2 = 6.375183$ depend on the lattice shape
and are related to the generalized zeta function of the momentum lattice.
Similar relations for nonzero momenta involve other phase shifts $\delta_l(k)$ as well.
Relation (11.88) can be generalized to particles with unequal mass and
nonzero total momentum. The leading term always is proportional to the
inverse volume [30].

\subsection{Beyond the threshold region}
Asymptotically only stable states can be observed. Thus in full QCD simulations resonances usually will have to be identified by their impact on intermediate states. Utilizing equations like (11.88), L\"uscher [26, 30--32] derived a
relation between the two-particle eigenstates energies on the finite lattice and
the scattering phase shift, valid below the inelastic threshold.
Let us discuss the idea in a simple situation. Consider the scattering of
two bosons in a $1{+}1$ dimensional system of finite spatial extension $L$, but
infinite time extension. This situation is most conveniently described by using
a wave function for the relative motion of the scattering partners. The basic
assumption is that the interaction range is finite and smaller than $L$. Thus
the wave function outside the interaction range is a plane wave. The effect of
the interaction is taken into account by a momentum-dependent phase shift
$\delta(k)$, acquired in the interaction region as illustrated in Fig.~11.2.

\begin{figure}[htbp]
  \centering
  \includegraphics[width=0.8\textwidth]{fig112.png}
  \caption*{Fig. 11.2. \emph{This figure illustrates the behavior of the wave
  function: Outside the interaction region it is an unperturbed plane wave which
  picks up an extra phase shift in the interaction region (indicated by the
  arrow)}.}
  \label{fig:112}
\end{figure}

Imposing periodic boundary conditions, the matching of the plane wave
with momentum $k$ at $x = L$ implies
\begin{equation}
e^{ikL + 2i\delta(k)} = e^{ik\cdot 0} = 1 .
\end{equation}
This gives rise to the quantization condition for the relative momenta $k_n$:
\begin{equation}
k_n L + 2\delta(k_n) = 2n\pi ,
\end{equation}
where $n$ is an integer. For vanishing $\delta$ one recovers the free case where the
momenta are $k_n = 2\pi n/L$. For the interacting situation nontrivial momenta
may be deduced from (11.90). If the functional form $\delta(k_n)$ were known, one
could use this relation to find the quantized values of the momentum in this
finite volume. On the other hand, given the momentum spectrum from some
measurement, (11.90) allows the determination of the phase shift $\delta(k_n)$ for
each $k_n$.
The momenta can be obtained from the energy values of the two particle states which are accessible in the simulation. For given $L$ one computes
the discrete levels $W_0, W_1, W_2, \ldots$ and from these the values of $k_n$ using the
dispersion relation
\begin{equation}
W_n = 2\sqrt{m^2 + k_n^2} .
\end{equation}
The technical problem lies in the precise determination of the single-particle
mass and of the energy levels $W_n$.
For the determination of the energy spectrum one has to use techniques
like the variational method (discussed in Sect.~6.3.3) considering correlation
functions of a sufficiently large number of interpolators with the correct quantum numbers, capable of representing the space of scattering states [33, 34],
including the coupled single-particle channels. Usually several of the lowest
energy eigenmodes can be determined with sufficient reliability. Varying the
spatial size $L$ of the system allows one to cover different values of the momentum. In [34] a simple 2D system was studied which couples a heavier and two
lighter bosons on the lattice, with mass and coupling parameters allowing for a
decay like in the $\rho \to \pi\pi$ system. Figure~11.3 demonstrates the expected phenomenon of level-crossing avoidance, which leads to a resonating phase shift.

\begin{figure}[htbp]
  \centering
  \includegraphics[width=0.9\textwidth]{fig113.png}
  \caption*{Fig. 11.3. \emph{Demonstration of the relation between scattering
  phase shifts and the finite volume energy spectrum in a simple toy model
  coupling a heavy particle to two light ones (figures from [34]). We show the
  energy levels observed without (a) and with (b) interaction coupling the single
  heavy particle to the two light ones. The horizontal dashed lines in (a) and (b)
  indicate the 2- and 4-particle thresholds, the horizontal full line in (a)
  denotes the single-(heavy)particle mass. Figure (c) shows the resulting
  resonance phase shift for (b); the vertical dashed line denotes the 4-particle
  threshold. For the 2D model the phase shift starts at $-\pi/2$ (no interaction).
  In all plots the symbols show the numerical data and the full curves represent
  models for the free (a) or resonating interacting case (b), (c)}.}
  \label{fig:113}
\end{figure}

One has to respect carefully the limitations of the approach: The interaction region and the single-particle correlation length ought to be smaller than
the spatial volume, in particular $mL \gg 1$. The relation is applicable only below the first inelastic threshold. Polarization effects due to virtual particles
running around the torus should be under control. Lattice artifacts will turn
up for large values of $k$.
In the physical 4D situation the relationship between phase shift, lattice
size, and momentum becomes somewhat more complicated:
\begin{equation}
\delta(k) = \phi(q) \bmod \pi , \qquad
q = \frac{kL}{2\pi} , \qquad
\tan(-\phi(q)) = \frac{q\,\pi^{3/2}}{\mathcal{Z}_{00}(1; q^2)} , \qquad
\phi(0) = 0 .
\end{equation}
Here $\mathcal{Z}_{00}$ denotes the Riemann zeta function. A 4D bosonic model has been
studied in [35, 36] and a quenched QCD simulation in [37].
The $\rho$-meson is the resonance in the $\pi\pi$ $p$-wave channel and therefore the
pions in the decay have nonzero momentum. In [38] (cf., also [39]) the relation
was extended to nonrest frame situations. This has the advantage that on the
same set of configurations both, zero and nonzero momentum cases, can be
studied simultaneously, resulting in more data points. Smaller volumes and
a smaller volume range can be utilized. A dynamical fermion study of the
$\rho$-decay has used this method [40, 41].
In [42] the decay is studied by analyzing the energy levels in a situation
when the $\rho$ mass lies below the two-pion mass. From the $N_T$-dependence the
transition amplitude $\rho \to \pi\pi$ was determined, and from this the effective
coupling.

\section{Matrix elements}
\label{sec:matrixelements}
When exploring the internal structure of a hadron one wants to know the spatial distribution of charge, matter, and spin. One also needs to know, however,
the momentum distribution. Both quantities are intrinsically nonperturbative
and therefore a challenge for lattice QCD.
The spatial distribution may be probed by elastic scattering processes,
giving rise to the concept of form factors and charge distributions. The momentum distribution can be explored experimentally in hard (high-energy)

\begin{figure}[htbp]
  \centering
  \includegraphics[width=0.8\textwidth]{fig114.png}
  \caption*{Fig. 11.4. \emph{Matrix elements of operators (squares) between
  mesonic states (circles)}.}
  \label{fig:114}
\end{figure}

collisions like deep-inelastic scattering or Drell--Yan processes. In the parton
model [43] quarks, antiquarks, and gluons are all partons and considered in
a hadron at very high momentum. The point-like probes identify the parton
probability distributions to find the partons in the infinite momentum frame
with some fraction $x$ of the hadron's longitudinal momentum and a momentum transfer $Q^2$.
A quark distribution function in a hadron $h$ is defined through a matrix
element in continuum light-cone notation
\begin{equation}
q(x) = \frac{1}{2p^+}\int \frac{d\lambda}{2\pi}\, e^{i\lambda x}\,
\ev{h(p)}{\bar{\Psi}(0)\gamma^+\Psi(\lambda n)}{h(p)} ,
\end{equation}
with $p^+ = (p^0 + p^3)/\sqrt{2}$, $n^\mu = (1, 0, 0, -1)$ and $n \cdot p = 1$. The quark field $\Psi$ is
connected from its position to infinity with a gauge transporter to make the
expression gauge-invariant.
The generalized parton distributions (GPDs) are an extension of this concept, introducing further variables like the spin and the transverse momentum
and correlating spatial and momentum values. Via this extension one hopes
to describe also scattering processes at lower energies [44, 45].
Wilson's operator product expansion (OPE) [46--49] relates the light-cone
matrix elements (11.93) or GPDs to matrix elements of local operators
$\ev{h}{O_i}{h'}$. One also needs renormalization factors $Z_{O_i}(a\mu)$ converting the lattice result to, e.g., the $\overline{\mathrm{MS}}$ scheme. The local operators can be simple quark
bilinears like $\bar{u}\gamma_\nu d$ (for form factors), they may include derivatives like $\bar{u}\gamma_\nu D_\mu d$
(for GPDs), or more quarks like $\bar{s}\gamma_\mu u\,\bar{u}\gamma_\mu d$ (for weak interaction processes). Figure~11.4 gives examples of the matrix elements that one has to
determine.

\subsection{Pion form factor}
Meson form factors are among the simplest matrix elements. As a prototype
example we discuss the electromagnetic form factor of the pion $F^{(\pi)}$ that
describes the coupling of the pion to the photon. It is defined by the matrix
element
\begin{equation}
\ev{\pi^+(p_f)}{V_\mu}{\pi^+(p_i)}_{\mathrm{cont}} = (p_f + p_i)_\mu\, F^{(\pi)}(Q^2) ,
\end{equation}
where $p_i$ and $p_f$ are the 4-momenta of the initial and final pions, $Q^2 = (p_f -
p_i)^2 \equiv -t$ is the space-like invariant momentum transfer squared, and
\begin{equation}
V_\mu = \frac{2}{3}\,\bar{u}\gamma_\mu u - \frac{1}{3}\,\bar{d}\gamma_\mu d
\end{equation}
is the isovector vector current normalized to unit electric charge at zero momentum transfer (i.e., the on-shell photon coupling): $F^{(\pi)}(0) = 1$. Its spacelike values are determined from experiment.
The connection to the continuum normalization gives another factor,
\begin{equation}
\ev{\pi^+(p_f)}{V_\mu}{\pi^+(p_i)}_{\mathrm{latt}}
= \frac{1}{\sqrt{2E_iE_f}}\,
\ev{\pi^+(p_f)}{V_\mu}{\pi^+(p_i)}_{\mathrm{cont}} .
\end{equation}
The pion form factor is an analytic function of $t$. Its space-like values can be
determined from the values on the boundary of its analyticity domain, i.e., the
cut along the positive real $t$-axis, the time-like region. Although the cut starts
at $t = 4m_e^2$, significant contributions only come from the hadronic resonance
region, starting with the two-pion threshold at $4M_\pi^2$. Due to unitarity along
the first part of the cut, until inelastic channels become important, the phase
shift is essentially the $p$-wave phase shift of the elastic two-pion channel. That
region is dominated by the $\rho$-meson resonance. This feature has been exploited
by various dispersion relation representations [50--53].
Although there are further inelastic contributions from the four-pion channel, these remain tiny until the $\pi\omega$ channel opens. There is also a small contribution from the isoscalar $\omega$ coupling through higher-order electromagnetic
interactions. All these contributions show small effects on the near space-like
region. The dominance of the vector meson motivates a simple pole parameterization for this region.
The mean charge radius squared $\ev{r^2}_V$ is defined through
\begin{equation}
F^{(\pi)}(Q^2) = 1 - \frac{1}{6}\,\ev{r^2}_V\, Q^2 + O(Q^4)
\quad\Rightarrow\quad
\ev{r^2}_V \equiv 6\,\frac{dF^{(\pi)}(t)}{dt}\Big|_{t=0} .
\end{equation}
The current PDG average for its value is $0.45(1)~\mathrm{fm}^2$ [2].
With present lattice tools one studies the pion form factor in the space-like
region. For this we need to evaluate off-forward matrix elements (preferably
at several transferred momenta), i.e., the expectation values of
\begin{equation}
\sum_{y;\,y_0=\tau}\;\sum_{x;\,x_0=t}
\tr\Bigl[e^{iq\cdot y}\, S(0, y)\, O(y)\,
e^{-ip\cdot x}\, S(y, x)\,\gamma_5\, S(x, 0)\,\gamma_5\Bigr] ,
\end{equation}
where $S(y, x)$ is the quark propagator from $x$ to $y$ and $O(y)$ denotes the
operator inserted at $y$. We use the notations $p \equiv p_f$, $r \equiv p_i$ and denote

\begin{figure}[htbp]
  \centering
  \includegraphics[width=0.8\textwidth]{fig115.png}
  \caption*{Fig. 11.5. \emph{Schematic diagram for the matrix element (11.98)}.}
  \label{fig:115}
\end{figure}

the momentum transfer as $q = p - r$ (cf. Fig.~11.5). Ideally the time distances
between source and operator insertion and between the operator insertion and
the sink are long enough, such that the higher excitations in the pion channel
may be disregarded and the pion ground state dominates the insertion region.
Smearing of source and sink operators also helps to get a better ground state
signal.
In evaluating matrix elements one has to compute expectation values of
combinations of quark propagators as in Figs.~11.4 and~11.5. Each calculation
of a quark propagator starting from a given source requires an inversion of the
Dirac operator and thus a significant amount of computer time. One therefore
wants to reduce the number of propagator calculations. This is achieved by
using the so-called sequential source method [54, 55]. We first calculate the
quark propagator from the source at the origin to all other points of the
lattice as discussed in Chap.~6. This provides us with both the propagator
from the source origin to $x$ and, due to the $\gamma_5$-hermiticity property, also from
$y$ to the origin. All we need is another propagator from $x$ to $y$. The matrix
element is then written as
\begin{equation}
\sum_{y;\,y_0=\tau}
\tr\Bigl[e^{iq\cdot y}\, S(0, y)\, O(y)\, \Sigma(y, 0)\, \gamma_5\Bigr] ,
\end{equation}
where the sequential propagator is defined as
\begin{equation}
\Sigma(y, 0) = \sum_{x;\,x_0=t} e^{-ip\cdot x}\, S(y, x)\,\gamma_5\, S(x, 0) .
\end{equation}
It can be easily computed by an additional inversion of the Dirac operator $D$
for each choice of the final momentum $p$:
\begin{equation}
\sum_y D(z, y)\, \Sigma(y, 0) = e^{-ip\cdot z}\, \gamma_5\, S(z, 0) , \qquad z_0 = t .
\end{equation}
Changing the properties of the sink requires the computation of new sequential
propagators, and so simulating several final momenta, different field interpolators, or a different smearing for the sink rapidly becomes rather expensive.
For this reason one often works only at one value of the final momentum.
This approach then allows not only the determination of the mentioned matrix element but, with no more effort, also the consideration of other operator
insertions, e.g., for the scalar form factor.
One extracts the physical matrix elements by computing ratios of 3-point
and 2-point correlators. To see this, let us look at the 3-point correlator for inand outgoing pseudoscalar interpolating field operators $P$ and an operator $O$
at insertion time $\tau$. At $\tau$ we insert a complete set of states. For large enough
separation distance this can be saturated by the ground state in the channel,
the pion, and we find
\begin{equation}
\begin{aligned}
\ev{0}{P(t; p)\, O(\tau)\, P(r)}{0}
&= \ev{0}{P(t; p)}{\pi(p)}\,
\frac{e^{-E_p(t-\tau)}}{2E_p}\\
&\qquad \ev{\pi(p)}{O(\tau)}{\pi(r)}\,
\frac{e^{-E_r\tau}}{2E_r}\,
\ev{\pi(r)}{P(0; r)}{0} ,
\end{aligned}
\end{equation}
and for the pseudoscalar propagator,
\begin{equation}
\ev{0}{P(t; p)\, P(p)}{0}
= \ev{0}{P(t; p)}{\pi(p)}\,
\frac{e^{-E_p t}}{2E_p}\,
\ev{\pi(p)}{P(0; p)}{0} .
\end{equation}
Here we neglect the lattice periodicity in time direction, otherwise there will
be further terms from the backward propagation. This has been discussed,
e.g., in [56].
One now constructs a suitable combination of these functions, such that
the unknown factors and the exponentials cancel:
\begin{equation}
\begin{aligned}
R(t, \tau; p, q \equiv p - r)
&= \sqrt{2E_iE_f}\,
\frac{\ev{P(t; p)\, O(\tau; q)\, P(0; r)}}{\ev{P(t; p)\, P(0; p)}} \\
&\times
\sqrt{
  \frac{\ev{P(t; p)\, P(0; p)}}{\ev{P(t; r)\, P(0; r)}}
  \frac{\ev{P(\tau; p)\, P(0; p)}}{\ev{P(\tau; r)\, P(0; r)}}
  \frac{\ev{P(t - \tau; r)\, P(0; r)}}{\ev{P(t - \tau; p)\, P(0; p)}}} \\
&= \ev{\pi(p)}{O(\tau; q)}{\pi(r)} .
\end{aligned}
\end{equation}
One can keep the sink fixed at time $t$ and vary the time slice $\tau$ where the
operator $O$ sits (scanning a range of time slices). As a consequence, $R(t, \tau; p, q)$
exhibits two plateaus in $\tau$: $0 \ll \tau \ll t$ and $t \ll \tau \ll N_T$. In the case in
which the sink is put at $t = N_T/2$ the ratio is (in our case) anti-symmetric in
$\tau - N_T/2$. The value of the form factor can be read off from the plateau values.
Further normalization factors due to smearing of source and sink operators
cancel also in the ratio.
The resulting plateau value for $Q^2 = 0$ will not necessarily be 1, which is
the electric charge in our example. This is because the naive lattice current
$V_\mu$ of (11.95) is not conserved. One could replace it by the conserved current,
which, however, is numerically demanding for chiral actions, as discussed in
Sect.~\ref{sec:naivecurrents}. For other actions even the conserved current will need $O(a)$
correction terms to obtain improved currents (see Sect.~9.1). Taking the simple
point-like current one has to determine the renormalization factor $Z_V$ relating
it to the renormalized current $V_\mu^{(r)}$.
Source, insertion, and sink are located on three time slices. The source
operator usually is localized on or around a site. In a Fourier decomposition
it therefore contributes to all possible 3-momenta. By projecting the sink to a
3-momentum $p$ and the insertion to $q$, the relevant component of the source is
$r = p - q$. All other contributions cancel in the sum over all field configurations
due to translation invariance. In a realistic situation with a large but finite
number of gauge configurations the cancellation is only approximate.
For the electromagnetic pion form factor a suitable choice of the 3momenta is $|p_f| = |p_i|$ such that $E_i = E_f$. In that situation $Q^2 = |\mathbf{q}|^2$
and combining (11.94) and (11.96) the kinematical factors simplify due to
\begin{equation}
\frac{E_i + E_f}{\sqrt{2E_iE_f}} = 1 .
\end{equation}
This removes another possible source of statistical errors. A possible way to
minimize the computational effort is to fix the momentum of the sink to
$p = (1, 1, 0)\,\Delta p$ for $\Delta p = 2\pi/(aN)$, the spatial momentum gap. Thus one has
to compute one further (sequential) propagator for each configuration. The 3point function is then evaluated for $q$ values such that $|r| = |p - q| = \sqrt{2} = |p|$.
These are the 12 combinations $(0, 0, 0)$, $(2, 2, 0)$, $(2, 0, 0)$, $(0, 2, 0)$, $(0, 1, \pm 1)$,
$(1, 0, \pm 1)$, $(2, 1, \pm 1)$, $(1, 2, \pm 1)$. This leads to results for $Q^2 = 2n(\Delta p)^2$ with
$n = 0, 1, 2, 3, 4$.
The momentum projection for the 2-point and 3-point functions gives rise
to significant statistical noise, in particular close to the symmetry point $t =
N_T/2$. One therefore chooses the position of the sink at time slices somewhat
closer to the source. The matrix elements are then obtained by combining the
plateau values on either side of the sink $R_{\mathrm{lhs}} \pm R_{\mathrm{rhs}}$ where the relative sign
depends on the parity properties of the interpolator. The sign is negative for
the vector form factor.
The earliest lattice studies of the pion form factor were [57, 58]. The spacelike form factor usually is well approximated by the time-like $\rho$ pole, which
in quenched calculations [59, 60] at unphysical large quark masses is larger
than the experimental value. For simulations with dynamical quarks, extrapolating to physical quark masses, agreement with the physical $\rho$ mass and
the measured pion form factor improves [56, 61, 62]. The pion is at the heart
of chiral symmetry breaking and therefore ChPT provides several constraints
[53, 63--65].

\subsection{Weak matrix elements}
Weak interactions couple almost point-like to quarks, but the quarks live
in the QCD-dominated environment of the hadron. Nonperturbative effects
define the hadron structure and therefore are important to analyze the way
the weak interactions ``see'' the hadron.
In the standard model (SM) the electroweak interactions are mediated by
the photon and the heavy vector bosons $W^\pm$ and $Z^0$. These bosons couple
to bilinear combinations of leptons and bilinear combinations of quarks, thus
allowing for quark flavor changing processes. The six flavors of quarks are
arranged in three families: $(u, d)$, $(c, s)$, and $(t, b)$. The weak interactions
``see'' the quark flavors slightly differently than the strong interactions. The
corresponding flavor eigenstates are denoted by $d'$, $s'$, and $b'$ and are related to
the mass eigenstates $d$, $s$, and $b$ via the Cabibbo--Kobayashi--Maskawa (CKM)
matrix [66, 67]
\begin{equation}
\begin{pmatrix} d' \\ s' \\ b' \end{pmatrix}
= V_{\mathrm{CKM}} \begin{pmatrix} d \\ s \\ b \end{pmatrix}
= \begin{pmatrix}
V_{ud} & V_{us} & V_{ub} \\
V_{cd} & V_{cs} & V_{cb} \\
V_{td} & V_{ts} & V_{tb}
\end{pmatrix}
\begin{pmatrix} d \\ s \\ b \end{pmatrix} .
\end{equation}
The other three quarks may be left unchanged without loss of generality. The
CKM matrix is complex (allowing for CP violation) and unitary (enforcing
absence of flavor changing neutral current transitions at tree level). For a
review on these aspects, see [68, 69].
With certain further phenomenologically motivated approximations the
matrix is often formulated in terms of four so-called generalized Wolfenstein
parameters: $\lambda = V_{us} = \sin\theta_c$, $A$, $\rho$, and $\eta$. These parameters may be related
to various experiments. A nonvanishing $\eta$ is responsible for CP violation in
the standard model.
Unitarity of the CKM matrix gives relations like
\begin{equation}
V_{ud}V_{ub}^* + V_{cd}V_{cb}^* + V_{td}V_{tb}^* = 0 .
\end{equation}
This relation can be represented as the ``unitarity triangle'' in the $(\rho, \eta)$ plane.
The endpoints of the triangle's base are $(0, 0)$ and $(1, 0)$ and the third corner
is then constrained by experiments, in particular by $K$ and $B$ physics results.
Lattice calculations may help in determining the QCD input relevant for the
computation of the CKM matrix elements [70--72].

\subsection{OPE expansion and effective weak Hamiltonian}
\label{sec:weak}
Weak decays of hadrons involve two mass scales: the heavy intermediate vector
bosons $W^\pm$ and $Z$ and the QCD low-energy scale dominating the hadron
structure. One separates the problem into distinct parts by introducing an
effective weak Hamiltonian
\begin{equation}
H_{\mathrm{eff}} = \frac{G_F}{\sqrt{2}} \sum_i V_{\mathrm{CKM}}\, C_i(\mu)\, Q_i(\mu) .
\end{equation}
Here $G_F$ is the Fermi constant and the $Q_i$ denote local operators, built from
quark and lepton fields. The CKM factors and the Wilson coefficients $C_i$ govern the weight of the operators. This is the so-called operator product expansion (OPE) [46--49]. One of the local operators is just the well-known product
of $(V - A)$ currents from the Fermi theory for $\beta$-decay. Expectation values
of these terms between hadronic states then are related to the experimental
observables.
The Wilson coefficients (short-distance scale) can be computed in perturbation theory (cf., the NLO calculation [68]). The long-distance (nonperturbative) scale is encoded in the $Q_i$ evaluated at a renormalization scale $\mu$.
Typically $\mu$ is of the order of 1--2~GeV or higher; it is, however, much smaller
than the weak interaction scale $M_W$ and therefore corrections
$O(\mu^2/M_W^2)$ are negligible. In lattice calculations one computes the hadronic matrix elements
$\ev{h}{Q_i}{h'}$, i.e., matrix elements of four-quark operators between hadronic
states (cf., Fig.~11.4).
When one uses lattice Dirac operators that do not obey the Ginsparg--
Wilson equation, the lack of chiral symmetry leads to significant complications [73]:
\begin{itemize}
\item For $K - \bar{K}$ mixing and computation of $B_K$ one needs operators like
$O_{LL} = \bigl(\bar{s}\gamma_\mu(1 - \gamma_5)d\bigr)^2$. If chiral symmetry is not guaranteed, upon
renormalization these operators mix with other operators of different chirality
like $O_{PP} = \bigl(\bar{s}\gamma_5 d\bigr)^2$. Whereas the matrix element $\ev{K}{O_{LL}}{\bar{K}}$ vanishes in
the chiral limit (proportional to the quark mass), the element $\ev{K}{O_{PP}}{\bar{K}}$
approaches a constant and therefore high accuracy is necessary to separate
the physically interesting part.
\item Operators of the effective Hamiltonian mix with lower dimensional operators. An example is the $\Delta S = 1$ Hamiltonian relevant to $K \to \pi\pi$. The
dimension-6 operator $\bar{s}\gamma_\mu(1 - \gamma_5)u\;\bar{u}\gamma_\mu(1 - \gamma_5)d$ mixes with operators like
$\bar{s}d$, $\bar{s}\gamma_5 d$, and $\bar{s}\gamma_\mu\gamma_\nu G_{\mu\nu} d$ (where $G$ denotes the gluon field). The mixing
coefficients diverge in powers of $1/a$ and precise nonperturbative methods
have to be employed to deal with these problems.
\end{itemize}

With actions like the Wilson action, that violate chiral symmetry, computations are therefore cumbersome due to the need to subtract divergent terms.
Using Ginsparg--Wilson fermions, on the other hand, improves that situation
significantly. For the $\Delta I = \tfrac{1}{2}$ problem one does not have any power divergent
subtractions [74] and for the computation of $\varepsilon'/\varepsilon$ subtraction is needed (for a
review see, e.g., [13]).
Nonleptonic weak decays, where the final state consists exclusively out of
hadrons, like $K \to \pi\pi$, are a challenging task. There are two main groups
of problems. The ultra-violet problem is to construct finite matrix elements
of renormalized operators from the bare lattice operators; this is discussed in
[6, 75--77]. The infrared problem has to do with re-scattering of the outgoing
two-particle system, the continuation from Euclidean to Minkowski space--
time and the effects of the spatial volume.
The Maiani--Testa theorem [78] states that in Euclidean calculations in
very large spatial volumes the matrix element is dominated by the pair of
zero momentum single-particle states. In infinite volume in Euclidean correlation functions one obtains the average of matrix elements into in- and out
two-pion states. This situation improves, however, when dealing with finite
volumes. Lellouch and L\"uscher [79] derived a relation between the $K \to \pi\pi$
matrix elements in finite volume and the physical kaon-decay amplitudes. In
accessible finite volumes the two-pion energy spectrum is far from being continuous and a kaon at rest cannot decay into two pions (unless one of the
two-particle energy levels is close to its mass). A simple expression then relates the transition amplitude in finite volume to the decay rates in infinite
volume. The work has been adapted [80] to include all elastic states below the
inelastic threshold. The methods allow to extract the decay amplitudes also
when the kaon mass is not equal to the two-pion energy, i.e., when an inserted
(weak Hamiltonian) operator carries nonvanishing energy momentum. For a
review on lattice results for kaon physics see [81] and, more generally, for what
lattice QCD can do for the standard model, see [82].

\section*{References}
\begin{enumerate}[label={[\arabic*]}]
  \item M. L\"uscher: in \emph{Probing the Standard Model of Particle Interactions}, Proceedings Les Houches 1997, edited by R. Gupta, A. Morel, E. de Rafael, and F. David, Vol. 1 (Elsevier, Amsterdam 1998) 268
  \item C. Amsler et al.\ [Particle Data Group]: Phys.\ Lett.\ B 667, 1 (2008) 270, 291
  \item M. Gell-Mann, R. J. Oakes, and B. Renner: Phys.\ Rev.\ 175, 2195 (1968) 274
  \item S. Weinberg: Physica A 96, 327 (1979) 274
  \item J. Gasser and H. Leutwyler: Nucl.\ Phys.\ B 250, 465 (1985) 274
  \item M. Bochicchio et al.: Nucl.\ Phys.\ B 262, 331 (1985) 274, 296
  \item P. Hasenfratz: Nucl.\ Phys.\ B 525, 401 (1998) 274
  \item P. Hasenfratz et al.: Nucl.\ Phys.\ B 643, 280 (2002) 275, 276, 277, 278
  \item Y. Kikukawa and A. Yamada: Nucl.\ Phys.\ B 547, 413 (1999) 276
  \item S. Capitani et al.: Phys.\ Lett.\ B 468, 150 (1999) 277
  \item L. Giusti, C. Hoelbling, and C. Rebbi: Phys.\ Rev.\ D 64, 114508 (2001) 277
  \item C. Gattringer, P. Huber, and C. B. Lang: Phys.\ Rev.\ D 72, 094510 (2005) 279
  \item S. Capitani: Phys.\ Rept.\ 382, 113 (2003) 280, 281, 284, 296
  \item G. Martinelli et al.: Nucl.\ Phys.\ B 445, 81 (1995) 281, 282, 283
  \item E. Franco and V. Lubicz: Nucl.\ Phys.\ B 531, 641 (1998) 281, 284
  \item M. L. Paciello, S. Petrarca, B. Taglienti, and A. Vladikas: Phys.\ Lett.\ B 341, 187 (1994) 282
  \item L. Giusti et al.: Int.\ J.\ Mod.\ Phys.\ A 16, 3487 (2001) 282, 283
  \item L. Giusti, S. Petrarca, B. Taglienti, and N. Tantalo: Phys.\ Lett.\ B 541, 350 (2002) 282
  \item M. G\"ockeler et al.: Nucl.\ Phys.\ B 544, 699 (1999) 282, 283, 284
  \item H. Suman and K. Schilling (unpublished) arXiv:hep-lat/9306018 (1993) 283
  \item J. E. Mandula and M. Ogilvie: Phys.\ Lett.\ B 48, 156 (1990) 283
  \item K. G. Chetyrkin and A. Retey: Nucl.\ Phys.\ B 583, 3 (2000) 284
  \item J. A. Gracey: Nucl.\ Phys.\ B 662, 247 (2003) 284
  \item V. Gimenez, L. Giusti, F. Rapuano, and M. Talevi: Nucl.\ Phys.\ B 531, 429 (1998) 284
  \item S. Capitani, M. L\"uscher, R. Sommer, and H. Wittig: Nucl.\ Phys.\ B 544, 669 (1999) 284
  \item M. L\"uscher: Commun.\ Math.\ Phys.\ 104, 177 (1986) 284, 286, 287
  \item P. Hasenfratz and H. Leutwyler: Nucl.\ Phys.\ B 343, 241 (1990) 286
  \item A. Hasenfratz et al.: Z.\ Physik C 46, 257 (1990) 286
  \item A. Hasenfratz et al.: Nucl.\ Phys.\ B 356, 332 (1991) 286
  \item M. L\"uscher: Commun.\ Math.\ Phys.\ 105, 153 (1986) 286, 287
  \item M. L\"uscher: Nucl.\ Phys.\ B 354, 531 (1991) 287
  \item M. L\"uscher: Nucl.\ Phys.\ B 364, 237 (1991) 287
  \item M. L\"uscher and U. Wolff: Nucl.\ Phys.\ B 339, 222 (1990) 288
  \item C. R. Gattringer and C. B. Lang: Nucl.\ Phys.\ B 391, 463 (1993) 288, 289
  \item J. Nishimura: Phys.\ Lett.\ B 294, 375 (1992) 289
  \item M. G\"ockeler, H. Kastrup, J. Westphalen, and F. Zimmermann: Nucl.\ Phys.\ B 425, 413 (1994) 289
  \item S. Aoki et al.: Phys.\ Rev.\ D 67, 014502 (2003) 289
  \item K. Rummukainen and S. Gottlieb: Nucl.\ Phys.\ B 450, 397 (1995) 289
  \item C. Kim, C. T. Sachrajda, and S. R. Sharpe: Nucl.\ Phys.\ B 727, 218 (2005) 289
  \item S. Aoki et al.: PoS LAT2006, 110 (2006) 289
  \item S. Aoki et al.: Phys.\ Rev.\ D 76, 094506 (2007) 289
  \item C. McNeile and C. Michael: Phys.\ Lett.\ B 556, 177 (2006) 289
  \item R. P. Feynman: \emph{Photon-Hadron Interactions} (Benjamin, New York 1972) 290
  \item D. Diakonov: Prog.\ Part.\ Nucl.\ Phys.\ 51, 173 (2003) 290
  \item X. Ji: Ann.\ Rev.\ Nucl.\ Part.\ Sci.\ 54, 413 (2004) 290
  \item K. G. Wilson: Phys.\ Rev.\ 179, 1499 (1969) 290, 295
  \item K. G. Wilson and W. Zimmermann: Commun.\ Math.\ Phys.\ 24, 87 (1972) 290, 295
  \item W. Zimmermann: Ann.\ Phys.\ 77, 570 (1973) 290, 295
  \item E. Witten: Nucl.\ Phys.\ B 120, 189 (1977) 290, 295
  \item J. Gounaris and J. J. Sakurai: Phys.\ Rev.\ Lett.\ 21, 244 (1968) 291
  \item M. F. Heyn and C. B. Lang: Z.\ Phys.\ C 7, 169 (1981) 291
  \item J. F. de Troc\'oniz and F. J. Yndur\'ain: Phys.\ Rev.\ D 65, 093001 (2002) 291
  \item H. Leutwyler: in \emph{Continuous Advances in QCD 2002}: Conference Proceedings Minneapolis, Minnesota, 2002, edited by K. A. Olive, M. A. Shifman, and M. B. Voloshin (World Scientific, Singapore 2002) 291, 294
  \item C. W. Bernard, T. Draper, G. Hockney, and A. Soni: in \emph{Lattice Gauge Theory: A Challenge in Large Scale Computing}, edited by B. Bunk, K. H. M\"utter, and K. Schilling (Plenum, New York 1986) 292
  \item G. W. Kilcup et al.: Phys.\ Lett.\ B 164, 347 (1985) 292
  \item D. Br\"ommel et al.: Eur.\ Phys.\ J.\ C 51, 335 (2007) 293, 294
  \item G. Martinelli and C. T. Sachrajda: Nucl.\ Phys.\ B 306, 865 (1988) 294
  \item T. Draper, R. Woloshyn, W. Wilcox, and K.-F. Liu: Nucl.\ Phys.\ B 318, 319 (1989) 294
  \item J. van der Heide, J. Koch, and E. Laermann: Phys.\ Rev.\ D 69, 094511 (2004) 294
  \item S. Capitani, C. Gattringer, and C. B. Lang: Phys.\ Rev.\ D 73, 034505 (2005) 294
  \item F. D. R. Bonnet et al.: Phys.\ Rev.\ D 72, 054506 (2005) 294
  \item S. Hashimoto et al.: PoS LAT2005, 336 (2005) 294
  \item H. Gausterer and C. B. Lang: Z.\ Phys.\ C 28, 475 (1985) 294
  \item G. Colangelo, J. Gasser, and H. Leutwyler: Nucl.\ Phys.\ B 603, 125 (2001) 294
  \item J. Bijnens: PoS LATTICE2007, 004 (2007) 294
  \item N. Cabibbo: Phys.\ Rev.\ Lett.\ 10, 531 (1963) 294
  \item M. Kobayashi and T. Maskawa: Prog.\ Theor.\ Phys.\ 49, 652 (1973) 294
  \item J. Buras: in \emph{Proceedings of the International School of Subnuclear Physics: 38th Course: Theory and Experiment Heading for New Physics}, Erice, Italy, 27 Aug--5 Sep 2000, edited by A. Zichichi (World Scientific, Singapore 2001) 295
  \item J. Buras: in \emph{Lectures given at the European CERN School}, Saint Feliu de Guixols, 2004 (CERN, Geneva, Switzerland 2005) 295
  \item Z. Ligeti: PoS LAT2005, 012 (2005) 295
  \item M. Okamoto: PoS LAT2005, 013 (2005) 295
  \item 5th International Workshop on the CKM Unitarity Triangle: \url{http://ckm2008.roma1.infn.it} (2008) 295
  \item A. Soni: Pramana 62, 415 (2004) 295
  \item S. Capitani and L. Giusti: Phys.\ Rev.\ D 64, 014506 (2001) 296
  \item L. Maiani, G. Martinelli, G. C. Rossi, and M. Testa: Phys.\ Lett.\ B 176, 445 (1986) 296
  \item L. Maiani, G. Martinelli, G. C. Rossi, and M. Testa: Nucl.\ Phys.\ B 289, 505 (1987) 296
  \item C. Bernard, T. Draper, G. Hockney, and A. Soni: Nucl.\ Phys.\ B (Proc.\ Suppl.) 4, 483 (1988) 296
  \item L. Maiani and M. Testa: Phys.\ Lett.\ B 245, 585 (1990) 296
  \item L. Lellouch and M. L\"uscher: Commun.\ Math.\ Phys.\ 219, 31 (2001) 296
  \item C.-J. Lin, G. Martinelli, C. Sachrajda, and M. Testa: Nucl.\ Phys.\ B 619, 467 (2001) 296
  \item L. Lellouch: PoS LATTICE2008, 015 (2009) 296
  \item T. DeGrand and C. DeTar: \emph{Lattice Methods for Quantum Chromodynamics} (World Scientific, Singapore 2006) 296
\end{enumerate}
